\documentclass[a4paper,12pt]{article}
\usepackage{../../mypackages}
\usepackage{../../macros}

\setlength{\parindent}{0pt}


\begin{document}

\title{Chapitre 2 : Les ions}
\author{N. Bancel}
\date{Octobre 2024}
\maketitle

\section{Connaissances à connaître pour le contrôle}

\begin{tcolorbox}[colback=blue!10, colframe=blue!80, title=]
  \begin{compactitem}
    \item Numéro atomique de l'atome
    \item Définition d'un ion (et la différence entre un cation et un anion)
    \item Les méthodes d'identification des ions
    \item Différence entre électriquement neutre et conducteur
  \end{compactitem}
\end{tcolorbox}

\section{Rappels : L'atome}

Activité : 5 minutes pour 
\begin{compactitem}
  \item Dessiner l’atome et identifier toutes ses particules et leurs caractéristiques (charge ?)
  
\begin{figure}[H]
  \centering
  \includegraphics[width=0.1\linewidth]{img/atome.jpg}
\end{figure}

  \item Dans le schéma ci-dessus, que représentent X, A, et Z ?
  \item Pourquoi peut-on dire qu’un atome est électriquement neutre ?
  \item Quelle est la différence entre un atome et une molécule ?
  \item Quelle est la formule brute de la molécule 
    \begin{compactitem}
      \item d'eau ?
      \item de dioxygène ?
      \item de dioxyde de carbone ?
      \item de dihydrogène ?
    \end{compactitem}
  \end{compactitem}




\section*{Les ions}

\begin{tcolorbox}[colback=green!10, colframe=green!80, title=Définition]
\textbf{Un ion est un atome (ou un groupe d’atomes) qui a perdu ou gagné un ou plusieurs électrons. Il est électriquement chargé.}

\par\vspace{1em}

Si des électrons sont perdus, l’ion formé est \textbf{positif} : c’est un \textbf{cation}.

\par\vspace{1em}

Si des électrons sont gagnés, l’ion formé est \textbf{négatif} : c’est un \textbf{anion}.
\end{tcolorbox}

\subsection*{Formule chimique}

La formule chimique d’un ion se compose du symbole chimique de l’atome (ou du groupe d’atomes) initial suivi de la charge de l’ion, inscrite en exposant.

\begin{figure}[H]
  \centering
  \includegraphics[width=0.7\linewidth]{img/al3+.png}
\end{figure}

\subsection*{Conséquences pour une solution}

Comme un ion est électriquement chargé, la présence d’ions rend une solution conductrice.

\begin{figure}[H]
  \centering
  \includegraphics[width=0.9\linewidth]{img2.jpg}
  \caption{\label{} Les ions}
\end{figure}

\subsection*{Trouver la formule d’un ion}

\begin{compactitem}
  \item En vous aidant de votre classification périodique, indiquez :
  \begin{compactitem}
    \item le nombre de \textbf{protons} contenus dans l’atome d’oxygène ;
    \item le nombre d’\textbf{électrons} contenus dans cet atome.
  \end{compactitem}

  \item En s’ionisant, l’atome d’oxygène \textbf{gagne deux électrons}.
  \begin{compactitem}
    \item Combien l’ion oxygène contient-il de \textbf{protons} et d’\textbf{électrons} ?
    \item Après avoir fait la \textbf{somme algébrique des charges} pour l’ion oxygène, donnez sa \textbf{formule}.
  \end{compactitem}
\end{compactitem}

\subsection*{L’ion bromure a pour formule $\mathrm{Br}^-$.}

\begin{compactitem}
  \item Combien de \textbf{protons} le noyau de l’ion bromure contient-il ?
  \item Faites la \textbf{somme algébrique des charges} pour cet ion.\\
  Combien d’\textbf{électrons} contient-il ?
\end{compactitem}

\subsection*{Trouver le nombre d'électrons}

\textbf{Certains pansements pour brûlures contiennent des ions argent, \ce{Ag^+}, qui ont des propriétés antibactériennes et cicatrisantes.}

\begin{compactitem}
  \item L’atome d’argent contient 47 protons. Combien l’ion argent possède-t-il de \textbf{protons} ?
  \item Faites la \textbf{somme algébrique des charges} pour l’ion argent, connaissant sa formule.\\
  Combien d’\textbf{électrons} cet ion contient-il ?
  \item Comment l’ion argent s’est-il formé à partir de l’atome d’argent ?\\
  (L’atome a-t-il \textbf{gagné} ou \textbf{perdu} un électron pour devenir un ion ?)
  \item Compléter :
  \begin{compactitem}
    \item Un ion est positif s’il a \dotfill un ou plusieurs \dotfill .
    \item Un ion est négatif s’il a \dotfill un ou plusieurs \dotfill .
  \end{compactitem}
\end{compactitem}

\subsection{Charge et formule de quelques ions}

\renewcommand{\arraystretch}{2.0} % augmente la hauteur des lignes
\begin{center}
\begin{tabular}{|l|c|c|c|c|}
\hline
\textbf{Nom} & \textbf{Nombre de protons} & \textbf{Nombre d’électrons} & \textbf{Charge de l’ion} & \textbf{Formule} \\
\hline
Ion fluorure & 9 & \hspace{2em} & $-1$ & \hspace{3em} \\
\hline
Ion zinc & 30 & \hspace{2em} & \hspace{2em} & \hspace{3em} \\
\hline
Ion fer II & 26 & \hspace{2em} & $+2$ & \hspace{3em} \\
\hline
Ion cuivre II & 29 & 27 & \hspace{2em} & \hspace{3em} \\
\hline
Ion \hspace{2em} & 13 & 10 & $+3$ & \hspace{3em} \\
\hline
Ion hydrogène & 1 & 0 & $+1$ & \hspace{3em} \\
\hline
\end{tabular}
\end{center}


\section*{Identification d'ions dans une solution}

\subsection*{Concept}

Une solution étant électriquement neutre, si elle contient des cations, elle contient nécessairement des anions aussi.

\begin{tcolorbox}[colback=purple!10, colframe=purple!80]
\textbf{Le nom donné à une solution renseigne sur sa composition.}
\end{tcolorbox}

Par exemple, une solution de chlorure de sodium contient des ions chlorure Cl\textsuperscript{-} et des ions sodium Na\textsuperscript{+}. Il est sous-entendu que le solvant est ici l'eau.

\begin{figure}[H]
  \centering
  \includegraphics[width=0.9\linewidth]{img4.jpg}
  \caption{\label{} Les ions}
\end{figure}

\subsection*{Identification d'ions}

\begin{tcolorbox}[colback=purple!10, colframe=purple!80]
\textbf{Il est possible d'identifier certains ions présents dans une solution grâce à des tests caractéristiques.} Au contact d'un réactif adapté, chaque ion forme un précipité caractéristique.
\end{tcolorbox}

Ces tests se réalisent sur un \textbf{échantillon} de la solution et jamais sur la totalité de la solution.

\begin{figure}[H]
  \centering
  \includegraphics[width=0.9\linewidth]{img3.jpg}
  \caption{\label{} Les ions}
\end{figure}

\begin{figure}[H]
  \centering
  \includegraphics[width=0.9\linewidth]{img/tests.jpg}
\end{figure}

\subsection*{Problème Brevet}

\begin{figure}[H]
  \centering
  \includegraphics[width=0.9\linewidth]{img/brevet_ions.jpg}
\end{figure}

\end{document}







